Zwischen zwei Stössen fliegt ein Teilchen geradlinig mit konstanter Geschwindigkeit. Die mittlere Zeit ist die mittlere freie Weglänge $\lambda$ geteilt durch die mittlere Schnelligkeit $\bar v = \sqrt{8kT/(\pi m)} = \sqrt{8RT/(\pi M)}$ (FoTaBe S.~171).

Für $\mathrm{N_2}$ ist $M = \MNzF = 2\cdot\MNaO$ (FoTaBe S.~228), und mit $T = \TSF = \TS$ ist
\begin{align}
 \bar v &= \vSF = \sqrt{\frac{8\cdot\ncR\cdot\TS}{\pi\cdot\MNz}} = \vS
\end{align}
Damit wird
\begin{align}
 \tau &= \tStF = \frac{\lam}{\vS} = \tSt \approx \result{\tStP{-9}{2}}
\end{align}
Ein Stickstoffmolekül stösst also rund zehn Milliarden Mal pro Sekunde mit anderen Molekülen zusammen.
